% Problem record op_aa21ac5ebca8b888. This TeX file is derived from
% database/problems_json/op_aa21ac5ebca8b888.json by scripts/sync-tex.mjs; edit the JSON record
% and rerun the script rather than editing this file.
\section{Entanglement cost of a qubit Bell-diagonal state}
\label{sec:op_aa21ac5ebca8b888}

\paragraph{Problem.}
What is the entanglement cost of a qubit Bell-diagonal state for an arbitrary
probability vector $\mathbf p=(p_I,p_X,p_Y,p_Z)$?  Let
$\lvert\Phi^+\rangle=(\lvert00\rangle+\lvert11\rangle)/\sqrt2$ and
$\lvert\Phi_P\rangle=(I\otimes P)\lvert\Phi^+\rangle$ for
$P\in\{I,X,Y,Z\}$.  The state is
\begin{equation}
  \rho_{\mathbf p}
  :=\sum_{P\in\{I,X,Y,Z\}}p_P
      \lvert\Phi_P\rangle\!\langle\Phi_P\rvert,
  \qquad p_P\geq0,
  \qquad \sum_{P\in\{I,X,Y,Z\}}p_P=1.
  \label{eq:p7-bell-diagonal-state}
\end{equation}
For the state in Eq.~\eqref{eq:p7-bell-diagonal-state}, determine
$E_C(\rho_{\mathbf p})$ for every $\mathbf p$ in the probability simplex,
where $E_C$ is the infimum asymptotic rate of ebits consumed by LOCC protocols
that prepare $\rho_{\mathbf p}^{\otimes n}$ with trace-norm error tending to
zero.

\subsection*{Status}

Unsolved

\subsection*{Source}

The problem is implicit in the regularized-entanglement-of-formation
characterization of entanglement cost and Wootters' one-copy formula for
two-qubit states \sourcecite{ref:p7-hayden-horodecki-terhal}{HHT01},
\sourcecite{ref:p7-wootters}{Woo98}.

\subsection*{Progress}

\begin{itemize}
  \item Let $p_\star:=\max\{p_I,p_X,p_Y,p_Z\}$,
  $C_{\mathbf p}:=\max\{0,2p_\star-1\}$, and
  $h_2(x):=-x\log_2x-(1-x)\log_2(1-x)$, with $0\log_2 0:=0$.  Regularized relative entropy of
  entanglement, entanglement cost, and Wootters' one-copy entanglement of
  formation give the rigorous bounds
  \begin{equation}
    L(p_\star)=E_R^\infty(\rho_{\mathbf p})
    \leq E_C(\rho_{\mathbf p})
    \leq E_F(\rho_{\mathbf p})
    =h_2\!\left(\frac{1+\sqrt{1-C_{\mathbf p}^{2}}}{2}\right),
    \qquad
    L(t):=
    \begin{cases}
      0, & 0\leq t\leq\tfrac12,\\
      1-h_2(t), & \tfrac12<t\leq1.
    \end{cases}
    \label{eq:p7-cost-bounds}
  \end{equation}
  Here $E_C=E_F^\infty$
  \sourcecite{ref:p7-hayden-horodecki-terhal}{HHT01}; Wootters gives the
  right endpoint \sourcecite{ref:p7-wootters}{Woo98}; and additivity of the
  relative entropy of entanglement gives
  $E_R^\infty(\rho_{\mathbf p})=L(p_\star)$
  \sourcecite{ref:p7-zhu-chen-hayashi}{ZCH10}.  Thus
  Eq.~\eqref{eq:p7-cost-bounds} solves $p_\star\leq1/2$ but generally leaves
  a gap for $p_\star>1/2$.  A recent semidefinite-programming construction
  gives a faithful, efficiently computable lower bound on
  $E_C(\rho_{\mathbf p})$ for every entangled two-qubit Bell-diagonal state,
  but does not close the gap
  \sourcecite{ref:p7-wang-jing-zhu}{WJZ25}.

  \item If at most two Bell probabilities are nonzero, say $q$ and
  $1-q$, then the entanglement of formation is strongly additive and
  \begin{equation}
    E_C(\rho_{\mathbf p})=E_F(\rho_{\mathbf p})
    =h_2\!\left(\frac12+\sqrt{q(1-q)}\right).
    \label{eq:p7-two-bell-cost}
  \end{equation}
  Vidal, D\"ur, and Cirac proved Eq.~\eqref{eq:p7-two-bell-cost} for a
  mixture of $\lvert\Phi^+\rangle$ and $\lvert\Phi^-\rangle$; local
  unitaries extend it to any pair of Bell states
  \sourcecite{ref:p7-vidal-dur-cirac}{VDC02}.  This settles the rank-at-most-two
  boundary of the family, not generic rank-three or rank-four states.

  \item On the isotropic line singled out in the question,
  $p_X=p_Y=p_Z$, write the probabilities and the known one-copy value as
  \begin{equation}
    p_I=F,
    \qquad p_X=p_Y=p_Z=\frac{1-F}{3},
    \qquad
    E_F(\rho_{\mathbf p})=
    \begin{cases}
      0, & 0\leq F\leq\tfrac12,\\
      h_2\!\left(\dfrac12+\sqrt{F(1-F)}\right),
        & \tfrac12<F\leq1.
    \end{cases}
    \label{eq:p7-isotropic-eof}
  \end{equation}
  Terhal and Vollbrecht determined this one-copy quantity
  \sourcecite{ref:p7-terhal-vollbrecht}{TV00}; it is also the qubit
  specialization of Wootters' formula.  For $1/2<F<1$, no proof is known
  that it equals the asymptotic LOCC entanglement cost.
\end{itemize}

\subsection*{References}

\begin{enumerate}[label={},leftmargin=4.8em,labelsep=0.5em]
  \footnotesize
  \item[\textup{[HHT01]}]\label{ref:p7-hayden-horodecki-terhal}
  P. M. Hayden, M. Horodecki, and B. M. Terhal,
  ``The Asymptotic Entanglement Cost of Preparing a Quantum State,''
  \emph{Journal of Physics A: Mathematical and General} \textbf{34},
  6891--6898 (2001).
  \href{https://doi.org/10.1088/0305-4470/34/35/314}{doi:10.1088/0305-4470/34/35/314};
  \href{https://arxiv.org/abs/quant-ph/0008134}{arXiv:quant-ph/0008134}.

  \item[\textup{[Woo98]}]\label{ref:p7-wootters}
  W. K. Wootters,
  ``Entanglement of Formation of an Arbitrary State of Two Qubits,''
  \emph{Physical Review Letters} \textbf{80}, 2245--2248 (1998).
  \href{https://doi.org/10.1103/PhysRevLett.80.2245}{doi:10.1103/PhysRevLett.80.2245};
  \href{https://arxiv.org/abs/quant-ph/9709029}{arXiv:quant-ph/9709029}.

  \item[\textup{[ZCH10]}]\label{ref:p7-zhu-chen-hayashi}
  H. Zhu, L. Chen, and M. Hayashi,
  ``Additivity and Non-Additivity of Multipartite Entanglement Measures,''
  \emph{New Journal of Physics} \textbf{12}, 083002 (2010).
  \href{https://doi.org/10.1088/1367-2630/12/8/083002}{doi:10.1088/1367-2630/12/8/083002};
  \href{https://arxiv.org/abs/1002.2511}{arXiv:1002.2511}.

  \item[\textup{[WJZ25]}]\label{ref:p7-wang-jing-zhu}
  X. Wang, M. Jing, and C. Zhu,
  ``Computable and Faithful Lower Bound on Entanglement Cost,''
  \emph{Physical Review Letters} \textbf{134}, 190202 (2025).
  \href{https://doi.org/10.1103/PhysRevLett.134.190202}{doi:10.1103/PhysRevLett.134.190202};
  \href{https://arxiv.org/abs/2311.10649}{arXiv:2311.10649}.

  \item[\textup{[VDC02]}]\label{ref:p7-vidal-dur-cirac}
  G. Vidal, W. D\"ur, and J. I. Cirac,
  ``Entanglement Cost of Bipartite Mixed States,''
  \emph{Physical Review Letters} \textbf{89}, 027901 (2002).
  \href{https://doi.org/10.1103/PhysRevLett.89.027901}{doi:10.1103/PhysRevLett.89.027901};
  \href{https://arxiv.org/abs/quant-ph/0112131}{arXiv:quant-ph/0112131}.

  \item[\textup{[TV00]}]\label{ref:p7-terhal-vollbrecht}
  B. M. Terhal and K. G. H. Vollbrecht,
  ``Entanglement of Formation for Isotropic States,''
  \emph{Physical Review Letters} \textbf{85}, 2625--2628 (2000).
  \href{https://doi.org/10.1103/PhysRevLett.85.2625}{doi:10.1103/PhysRevLett.85.2625};
  \href{https://arxiv.org/abs/quant-ph/0005062}{arXiv:quant-ph/0005062}.
\end{enumerate}

\subsection*{Comment}

Equation~\eqref{eq:p7-isotropic-eof} determines $E_C$ only for $F\leq1/2$
and $F=1$; for $1/2<F<1$ it is an $E_F$ formula.  The open task is to
regularize $E_F$ for entangled rank-three and rank-four Bell-diagonal states.
Problem~8 asks the analogous question for a different two-qubit family.

\subsection*{Field}

Entanglement theory

\subsection*{Topic}

Bell-diagonal states; Entanglement cost; Entanglement measures; Local operations and classical communication; Qubit systems

\subsection*{ID}

\texttt{op\_aa21ac5ebca8b888}
